%----------------------------------------------------------------------------------------
%	PHẦN 6: LỜI CẢM ƠN (ACKNOWLEDGEMENTS)
%----------------------------------------------------------------------------------------

\begin{acknowledgements}
	\addchaptertocentry{\acknowledgementname}
	\thispagestyle{empty}
	
	Lời đầu tiên, nhóm chúng em xin gửi lời cảm ơn chân thành và sâu sắc nhất đến các thầy cô giáo Trường Đại học Khoa học Tự nhiên, Đại học Quốc gia Hà Nội. Trong suốt quá trình học tập dưới mái trường, các thầy cô đã tận tình truyền dạy và trang bị cho nhóm những kiến thức nền tảng vô cùng quý báu về trí tuệ nhân tạo, đại số tuyến tính, giải tích,..., tạo tiền đề vững chắc để em thực hiện đề tài này.
	
	Đặc biệt, nhóm xin bày tỏ lòng biết ơn sâu sắc đến TS. Nguyễn Tiến Cường và CN. Vi Anh Quân. Sự hỗ trợ chuyên môn và những kiến thức bổ trợ quý báu về lý thuyết Học máy của hai thầy chính là kim chỉ nam giúp nhóm chúng em giải quyết bài toán “Xây dựng hệ thống nhận diện đa nhãn cảm xúc”, từ đó hoàn thiện nghiên cứu một cách bài bản, khoa học và hiệu quả hơn.
	
	Nhóm cũng xin cảm ơn các bạn cùng lớp Khóa 68 Kỹ thuật Điện tử và Tin học đã luôn chia sẻ tài liệu, đóng góp ý kiến phản biện và cổ vũ tinh thần cho nhóm trong suốt thời gian triển khai đề tài cuối kỳ môn Học Máy.
	
	Mặc dù đã có nhiều cố gắng để hoàn thiện đề tài một cách chỉn chu nhất, song do kiến thức và kinh nghiệm thực tế của bản thân còn hạn chế, quyển báo cáo chắc chắn không tránh khỏi những thiếu sót trong quá trình trình bày và thiết kế. Nhóm rất mong nhận được những ý kiến đóng góp và nhận xét từ quý thầy cô để có thể tiếp tục hoàn thiện đề tài cũng như nâng cao kiến thức và kỹ năng của bản thân trong thời gian tới.
	
	\begin{flushright}
		\begin{tabular}{r}
			Xin chân thành cảm ơn! \\[0.6cm]
			\textit{TP. Hà Nội, ngày 03 tháng 06 năm 2026} \\[0.2cm]
			\textbf{Nhóm sinh viên thực hiện} \\[0.6cm]
			Nhóm 6
		\end{tabular}
	\end{flushright}
\end{acknowledgements}